\input{./._relpath-to-latexroot.ltx}
\documentclass[\latexroot/\projectname]{subfiles}
\input{\latexroot/@resources/subfile-setup.ltx}

\begin{document}

\section{Motivation: two stylized facts}
\whenintegrated{\label{sec:motivation}} % integrated doc: SST for labels

Two patterns in the cross-country data on advanced-economy sovereign
debt motivate the framework that follows. First, a sharp \emph{kink}
at high debt: once total government debt crosses roughly $100\%$ of
GDP, banking-sector holdings rise much more steeply than below the
threshold. Second, substantial \emph{cross-country heterogeneity} at
any given debt level: countries with similar debt-to-GDP ratios place
their debt in very different hands.

\begin{figure}[htbp]
\centering
\begin{minipage}[t]{0.48\textwidth}
  \centering
  \includegraphics[width=\linewidth]{Figures/db_total_vs_totalD.png}
  \subcaption*{(a) Bank holdings vs.\ total debt.}
\end{minipage}\hfill
\begin{minipage}[t]{0.48\textwidth}
  \centering
  \includegraphics[width=\linewidth]{Figures/debt_composition_examples.png}
  \subcaption*{(b) Creditor composition at similar debt levels.}
\end{minipage}
\caption{Motivating evidence.}
\label{fig:motivation}
\end{figure}

%-------------------------------------------------------------
\subsection{Fact 1 — Banking takeover above a high-debt threshold}
%-------------------------------------------------------------

Once total government debt rises beyond roughly $100\%$ of GDP, the
banking sector begins to absorb a noticeably larger share of marginal
debt issuance than at lower debt levels. Below the threshold, an
additional unit of total debt translates into only a modest increase in
combined central-bank and private-bank holdings; above it, the same
incremental unit of debt is absorbed mostly by the banking system. The
slope of bank holdings against total debt steepens sharply around this
point, and banking eventually accounts for nearly all of the marginal
issuance. This is the empirical pattern we will refer to as the
``banking takeover'' and which the model in
Section~\ref{sec:model} is designed to rationalize.

%-------------------------------------------------------------
\subsection{Fact 2 — Cross-country heterogeneity at the same debt level}
%-------------------------------------------------------------

At any given debt-to-GDP ratio, advanced economies place their debt
with strikingly different creditor sectors. Some lean heavily on
domestic non-bank investors such as pension funds and insurance
companies; others rely on foreign investors; others place a large
share of their debt with the domestic banking system. This
heterogeneity does not vanish once we condition on the level of debt:
countries with intrinsically larger non-bank financial sectors absorb
a larger share of government debt through non-banks, while countries
with structurally smaller non-bank sectors push more debt onto
foreigners and the banking system.

%-------------------------------------------------------------
\subsection{What this implies for a framework}
%-------------------------------------------------------------

Together the two facts pin down what a useful framework needs to
deliver. It must (i) generate a sharp threshold in marginal banking
absorption as total debt rises, and (ii) allow the cross-country mix
of creditors at a given debt level to vary with the country's domestic
financial structure. Section~\ref{sec:model} develops a
cost-minimization model that produces both features: domestic
non-banks have country-specific capacity, foreign investors face a
common exposure ceiling, and banking is the residual margin with no
hard ceiling but a convex placement cost. Section~\ref{sec:empirical}
takes that framework directly to panel data and shows that both
stylized facts emerge as predicted.

% Smart bibliography: Only include bibliography if standalone AND has citations
\smartbib

\end{document}
